%!TEX root = ../../main.tex

\section{Convex Reformulations:\ Proofs}\label{app:convex-forms}

\begin{restatable}{lemma}{miFormulation}\label{lemma:mi-reformulation}
    The non-convex problem NC-ReLU (Problem \ref{convex-forms:eq:non-convex-relu-mlp}) is equivalent to the mixed-integer program,
    \begin{equation}\label{convex-forms:eq:mi-reformulation}
        \begin{aligned}
            \min_{W^{(1)}, W^{(2)}} \; & \calL\Big(\sum_{i = 1}^m \rbr{X W^{*(1)}_{\cdot, i}}_+ W^{*(2)}_{i}, y\Big) + \lambda \sum_{i = 1}^m \norm{W^{*(1)}_{\cdot, i}}_2 \\
                                       & \text{\emph{s.t.}} \; W^{(2)} \in \cbr{-1, 1}^m.
        \end{aligned}
    \end{equation}
\end{restatable}

\begin{proof}
    The proof proceeds in two steps: first we transform the objective into an equivalent problem which is invariant to certain scale re-parameterizations of the network parameters.
    Then, we use these scale re-parameterizations reduce the two optimization problems to each-other.

    Let \( p^* \) be the optimal value of the non-convex optimization problem and \( d^* \) the optimal value of the mixed-integer program.
    The ReLU activation function,
    \begin{align*}
        \rbr{a}_+ = \max\cbr{a, 0},
    \end{align*}
    is positively homogeneous, meaning \( \rbr{a * \beta}_+ = \beta \rbr{a}_+ \) for any scalar \( \beta \geq 0 \).
    Defining \( \bar{W}^{(1)}_{\cdot, i} = \beta_i W^{*(1)}_{\cdot, i}, \, \bar{W}^{(2)}_{i} = W^{*(2)}_{i} / \beta_i \), \( i \in [m] \), we have
    \begin{align*}
        f_{W^{(1)}, W^{(2)}}(X)
         & = \sum_{i=1}^m \rbr{X W^{*(1)}_{\cdot, i} }_+ W^{*(2)}_{i}
        = \sum_{i=1}^m \rbr{X W^{*(1)}_{\cdot, i} \frac{\beta_i}{\beta_i}}_+ W^{*(2)}_{i} \\
         & = \sum_{i=1}^m \rbr{X \bar{W}^{(1)}_{\cdot, i}}_+ \bar{W}^{(2)}_{i}
        = h(\bar{W}^{(1)}, \bar{W}^{(2)}),
    \end{align*}
    implying that the loss \( \calL(\sum_{i=1}^m \rbr{X W^{*(1)}_{\cdot, i} }_+
    W^{*(2)}_{i}, y) \) is invariant to ``scale-shift'' re-parameterizations of
    this form.
    To extend the invariance to the full objective function, recall Young's inequality,
    \[
        2 \abr{a, b} \leq a^2 + b^2,
    \]
    which yields,
    \begin{align*}
         & \calL\rbr{\sum_{i = 1}^m \rbr{X W^{*(1)}_{\cdot, i} }_+ W^{*(2)}_{i}, y } + \frac{\lambda}{2} \sum_{i = 1}^m \norm{W^{*(1)}_{\cdot, i}}_2 + \abs{W^{*(2)}_{i}}^2    \\
         & \hspace{3em} \geq \calL\rbr{ \sum_{i = 1}^m \rbr{X W^{*(1)}_{\cdot, i} }_+ W^{*(2)}_{i}, y } + \lambda \sum_{i = 1}^m \norm{W^{*(1)}_{\cdot, i}}\abs{W^{*(2)}_{i}}.
    \end{align*}
    For any choice of parameters \( W^{*(1)}_{\cdot, i}, W^{*(2)}_{i} \),
    equality in this expression is achieved with the rescaling
    \[
        \bar{W}^{(1)}_{\cdot, i} = W^{*(1)}_{\cdot, i} * \beta_i, \quad \quad \bar{W}^{(2)}_{i} = W^{*(2)}_{i} / \beta_i,
    \]
    where \( \beta_i = \sqrt{\frac{W^{*(2)}_{i}}{\norm{W^{*(1)}_{\cdot, i}}_2}} \).
    As this rescaling does not affect \( f_{W^{(1)}, W^{(2)}} \), it must be that any
    global minimizer \( \theta^* = \rbr{W^{*(1)}, W^{*(2)}} \) of
    Problem~\ref{convex-forms:eq:non-convex-relu-mlp} achieves the lower-bound
    in Young's inequality and,
    \begin{align}\label{convex-forms:eq:inner-product-reformulation}
         & \calL\rbr{\sum_{i = 1}^m \rbr{X W^{*(1)}_{\cdot, i}}_+ W^{*(2)}_{i}, y} + \frac{\lambda}{2} \sum_{i = 1}^m \norm{W^{*(1)}_{\cdot, i}}_2^2 + \abs{W^{*(2)}_{i}}^2 \\
         & \hspace{3em} =
        \calL\rbr{\sum_{i = 1}^m \rbr{X W^{*(1)}_{\cdot, i}}_+ W^{*(2)}_{i}, y} + \lambda \sum_{i = 1}^m \norm{W^{*(1)}_{\cdot, i}}_2\abs{W^{*(2)}_{i}}.
    \end{align}
    The right-hand side of this equation is invariant to scale
    re-parameterizations of the form
    \( \bar{W}^{(1)}_{\cdot, i} = \beta
    W^{*(1)}_{\cdot, i}, \bar{W}^{(2)}_{i} = W^{*(2)}_{i} / \beta \)
    for \( \beta > 0 \).
    Taking \( \beta = \abs{W^{*(2)}_{i}} \), we deduce
    \begin{align*}
        p^* = \calL\rbr{\sum_{i = 1}^m \rbr{X \bar{W}^{(1)}_{\cdot, i}}_+ \bar{W}^{(2)}_{i}, y} + \lambda \sum_{i = 1}^m \norm{\bar{W}^{(1)}_{\cdot, i}}_2 \geq d^*.
    \end{align*}

    To show the reverse inequality, observe that every global minimum \(
    \rbr{W^{*(1)}, W^{*(2)}} \) of
    Problem~\ref{convex-forms:eq:mi-reformulation} is trivially in the domain
    of the non-convex ReLU training problem.
    Using the mapping
    \[ \rbr{\bar{W}^{(1)}_{\cdot, i}, \bar{W}^{(2)}_{i}}
        = \rbr{\frac{W^{*(1)}_{\cdot, i}}{\sqrt{\norm{W^{*(1)}_{\cdot, i}}_2}},
            W^{*(2)}_i \sqrt{\norm{W^{*(1)}_{\cdot, i}}_2}},
    \]
    and plugging \( \rbr{\bar{W}^{(1)}, \bar{W}^{(2)}} \) into
    Problem~\ref{convex-forms:eq:non-convex-relu-mlp} shows \( d^* \geq p^* \).
    We have shown \( p^* = d^* \) and so the problems are formally equivalent
    with mappings between the solutions as given above.
\end{proof}

\convexDualityFree*
\begin{proof}
    The proof proceeds by showing the equivalence of C-ReLU and the mixed integer problem given in \cref{convex-forms:eq:mi-reformulation} and the invoking \cref{lemma:mi-reformulation}.
    Let \( p^* \) be the optimal value of the mixed-integer problem in~\eqref{convex-forms:eq:mi-reformulation} and \( d^* \) the optimal value of the convex program in~\eqref{convex-forms:eq:convex-relu-mlp}.
    We first show that \( d^* \geq p^* \).

    Suppose \( \rbr{v^*, u^*} \) is a global minimizer of Problem~\ref{convex-forms:eq:convex-relu-mlp} and let
    \begin{align*}
        \cbr{\rbr{W^{*(1)}_{\cdot, k}, W^{*(2)}_{k}}} = \bigcup_{D_i \in \tilde \calD_i} \cbr{\rbr{v^*_i, 1} : v^*_i \neq 0} \cup \cbr{\rbr{u^*_i, -1} : u^*_i \neq 0},
    \end{align*}
    where we set \( W^{*(1)}_{\cdot, k} = 0, \) and \( W^{*(2)}_{k} = 0 \) for all \( k \in [m], k > b \).
    It holds by assumption that \( b \leq m \) and thus \( \rbr{W^{*(1)}, W^{*(2)}} \) is a valid input for the mixed-integer problem.

    Recalling the constraints \( (2 D_i - I) X v^*_i \geq 0 \), and \( (2 D_i - I) X u^*_i \geq 0 \),
    we see that \( D_i X v_i^* = \rbr{X v_i^*}_+ \), \( D_i X u_i^* = \rbr{X u_i^*}_+ \), and thus
    \begin{align*}
        \rbr{X W^{*(1)}_{\cdot, k}} W^{(2)}_{j}
         & = \begin{cases}
                 D_i X v_i^*   & \mbox{if \( W^{*(1)}_{\cdot, k} = v_i^* \) for some \( i \in [b] \)} \\
                 - D_i X u_i^* & \mbox{if \( W^{*(1)}_{\cdot, k} = u_i^* \) for some \( i \in [b] \)} \\
                 0             & \mbox{otherwise}.
             \end{cases}
    \end{align*}
    Using this fact in the optimization objective for the convex program, we find
    \begin{align*}
        d^*
         & = \calL\rbr{\sum_{D_i \in \tilde \calD} D_i X (v_i - u_i), y} + \lambda \sum_{D_i \in \calD} \norm{v_i}_2 + \norm{u_i}_2         \\
         & = \calL\rbr{\sum_{k = 1}^m \rbr{X W^{*(1)}_{\cdot, k}}_+ W^{*(2)}_{k} - y} + \lambda \sum_{k = 1}^m \norm{W^{*(1)}_{\cdot, k}}_2 \\
         & \geq p^*,
    \end{align*}
    as required.

    To show the reverse inequality, let \( \rbr{W^{*(1)}, W^{*(2)}} \) be a
    solution to~\eqref{convex-forms:eq:mi-reformulation} and consider the
    set-function
    \[
        T(j) = \cbr{ i \in [m] : \text{diag}\rbr{X W^{*(1)}_{\cdot, i} > 0} = D_j }.
    \]
    Recalling
    \( \cbr{\text{diag}\rbr{X W^{*(1)}_{\cdot, i} > 0 : i \in [m]} } \subseteq \tilde \calD, \)
    by assumption, we define a valid candidate solution as
    \begin{equation*}
        \cbr{\rbr{v_j^*, u_j^*}}_{D_j \in \tilde \calD} = \cbr{ \sum_{i \in T(j)} W^{*(1)}_{\cdot, i} \mathds{1}(W^{*(2)}_{i} = 1), \sum_{i \in T(j)} W^{*(1)}_{\cdot, i} \mathds{1}(W^{*(2)}_{i} = -1) }_{D_j \in \tilde \calD}
    \end{equation*}

    We start by showing that the neurons indexed by \( T(j) \) can be merged without changing the objective of the mixed-integer problem.
    In particular, let \( j \in [m] \) be arbitrary and suppose that there
    exists \( l,k \in T(j) \) such that \( W^{*(2)}_{l} = W^{*(2)}_{k} = 1 \).
    By definition of \( T \), it holds that
    \( \rbr{X W^{*(1)}_{\cdot, l}}_+ \) and \( \rbr{X W^{*(1)}_{\cdot, k}}_+ \) have the same activation pattern.
    Accordingly, we have \(  \rbr{X W^{*(1)}_{\cdot, l}}_+ + \rbr{X W^{*(1)}_{\cdot, k}}_+ = \rbr{X (W^{*(1)}_{\cdot, l} + W^{*(1)}_{\cdot, k})}_+ \) by definition of the ReLU activation.
    Thus, merging these two parameter vectors as \( Z_{lk}^* = W^{*(1)}_{\cdot, l} + W^{*(1)}_{\cdot, k}\) does not change the prediction of the model in the mixed-integer program.

    Now we consider the group \( \ell_1 \) penalty term.
    Triangle inequality implies
    \begin{align*}
        \norm{\tilde Z_{lk}^*}_2 \leq \norm{W^{*(1)}_{\cdot, l}}_2 + \norm{W^{*(1)}_{\cdot, k}}_2,
    \end{align*}
    with equality if and only if \( W^{*(1)}_{\cdot, l} = 0 \), \( W^{*(1)}_{\cdot, k} = 0 \), or the vectors are collinear.
    Suppose that equality does not hold.
    Then the penalty term could be reduced setting \( W^{*(1)}_{\cdot, l} = Z_{lk} \) and \( W^{*(1)}_{\cdot, k} = 0\) while leaving the squared-loss term unchanged.
    But, this contradicts global optimality of \(W_{1}^*, W^{*(2)} \).
    Thus, it must be that \( W^{*(1)}_{\cdot, l} = 0 \), \( W^{*(1)}_{\cdot, k} = 0 \), or the vectors are collinear.
    In each case, we have that the merged vector \( \tilde Z_{lk}^* \) also attains the optimal value \( p^* \).
    Clearly a symmetric argument holds in the case \( W^{(2)}_{k} = W^{(2)}_{l} = -1 \).

    Arguing by induction if necessary, we deduce that the vectors \( v^*, u^* \) given by the solution mapping also attain \( p^* \).
    Recalling that \( D_j X v_j^* = \rbr{X v_j}_+ \) and \( D_j X u_j^* = \rbr{X u_j}_+ \) by choice of \( T(j) \) and definition of \( D_j \) gives
    \begin{align*}
        p^*
         & = \calL\rbr{\sum_{i=j}^P D_j X (v^*_j - u^*_j) - y} + \lambda \sum_{j=1}^P \norm{v^*_j}_2 + \norm{u^*_j}_2 \\
         & \geq d^*,
    \end{align*}
    where \( v_j^*, u_j^* \) are feasible.
    This completes the proof.
\end{proof}

\unconstrainedEquivalence*
\begin{proof}
    The proof proceeds similarly to the proof \cref{lemma:mi-reformulation}.

    Let \( p^* \) be the optimal value of the Problem~\ref{convex-forms:eq:gated-relu-mlp} and let \( d^* \) be the optimal value of the C-GReLU problem~\eqref{convex-forms:eq:unconstrained-relu-mlp}.
    For each \( z_i \in \tilde \calZ \) and any \( v \in \R^d \), we have the following equality by construction:
    \[ \phi_{z_i}(X, v) = \rbr{X z_i > 0} \circ X v = D_i X v. \]
    Accordingly, the non-convex optimization problem~\eqref{convex-forms:eq:gated-relu-mlp} can be written as
    \begin{equation}
        \min_{W^{(1)},W^{(2)}} \mathcal{L}\Big(\sum_{D_i \in \tilde \calD } D_i X W^{*(1)}_{\cdot, i} W^{*(2)}_{i}, y\Big) + \frac{\lambda}{2} \sum_{z_i \in \tilde \calZ} \norm{W^{*(1)}_{\cdot, i}}_2^2 + W^{*(2)}_{i},
    \end{equation}
    which makes the connection to C-GReLU clear.
    Applying Young's inequality gives,
    \begin{align*}
         & \mathcal{L}\Big(\sum_{D_i \in \tilde \calD } D_i X W^{*(1)}_{\cdot, i} W^{*(2)}_{i}, y\Big) + \frac{\lambda}{2} \sum_{z_i \in \tilde \calZ} \norm{W^{*(1)}_{\cdot, i}}_2^2 + W^{*(2)}_{i}           \\
         & \hspace{4em} \geq \mathcal{L}\Big(\sum_{D_i \in \tilde \calD } D_i X W^{*(1)}_{\cdot, i} W^{*(2)}_{i}, y\Big) + \lambda \sum_{z_i \in \tilde \calZ} \norm{W^{*(1)}_{\cdot, i}}_2 \abs{W^{*(2)}_{i}}
    \end{align*}
    For any choice of parameters \( W^{*(1)}_{\cdot, i}, W^{*(2)}_{i} \), equality in this expression is achieved with the rescaling
    \[
        \bar{W}^{(1)}_{\cdot, i} = W^{*(1)}_{\cdot, i} * \beta_i, \quad \quad \bar{W}^{(2)}_{i} = W^{*(2)}_{i} / \beta_i,
    \]
    where \( \beta_i = \sqrt{\frac{W^{*(2)}_{i}}{\norm{W^{*(1)}_{\cdot, i}}_2}} \).
    As this rescaling does not affect \( D_i X W^{*(1)}_{\cdot, i} \, W^{*(2)}_{i} \) for each \( i \), it must be that any global minimizer \( \theta^* = \cbr{W^{*(1)}_{\cdot, i}, W^{*(2)}_{i}} \) of Problem~\ref{convex-forms:eq:gated-relu-mlp} achieves the lower-bound in Young's inequality.
    Defining \( v_i' = W^{*(1)}_{\cdot, i} * W^{*(2)}_{i} \), we have shown
    \begin{align*}
        p^*
         & = \mathcal{L}\Big(\sum_{D_i \in \tilde \calD } D_i X W^{*(1)}_{\cdot, i} W^{*(2)}_{i}, y\Big) + \lambda \sum_{z_i \in \tilde \calZ} \norm{W^{*(1)}_{\cdot, i}}_2 \abs{W^{*(2)}_{i}} \\
         & = \mathcal{L}\Big(\sum_{D_i \in \tilde \calD } D_i X v_i', y \big) + \lambda \sum_{z_i \in \tilde \calZ} \norm{v_{i}'}_2
        \geq d^*,
    \end{align*}
    where we have used absolute homogeneity of the norm.

    To obtain the reverse inequality, let \( v^* \) be a global minimizer of C-GReLU and define \( \bar{W}^{(1)}_{\cdot, i} = \frac{v_i^*}{\sqrt{\norm{v_i^*}_2}} \), \( \bar{W}^{(2)}_{i} = \sqrt{\norm{v_i^*}_2} \) for all \( i \) to obtain

    \begin{align*}
        d^*
         & = \mathcal{L}\Big(\sum_{D_i \in \tilde \calD} D_i X v_i, y\Big)+ \lambda \sum_{D_i \in \tilde \calD} \norm{v_i}_2                                                                                                       \\
         & = \mathcal{L}\Big(\sum_{D_i \in \tilde \calD } D_i X \bar{W}^{(1)}_{\cdot, i} \bar{W}^{(2)}_{i}, y\Big) + \frac{\lambda}{2} \sum_{z_i \in \tilde \calZ} \norm{\bar{W}^{(1)}_{\cdot, i}}^2_2 + \abs{\bar{W}^{(2)}_{i}}^2
        \geq p^*,
    \end{align*}
    which completes the proof.

\end{proof}

%% Dimensionality reduction for unconstrained Problem

\begin{proposition}\label{prop:unconstrained-relu-mlp}
    Problem~\ref{convex-forms:eq:unconstrained-relu-mlp} is equivalent to following unconstrained relaxation of the ReLU training problem's convex reformulation (Problem~\ref{convex-forms:eq:convex-relu-mlp}):
    \begin{equation} \label{convex-forms:eq:unreduced-unconstrained-relu-mlp}
        \begin{aligned}
            \min_{v, u} \;
             & \half \calL\rbr{\sum_{D_i \in \tilde \calD} D_i X (v_i - u_i), y} + \lambda \sum_{D_i \in \tilde \calD} \norm{v_i}_2 + \norm{u_i}_2 \\
             & \text{s.t.} \;  \rbr{2 D_i - I_n} X v_i \succeq 0, \rbr{2 D_i - I_n} X u_i \succeq 0,
        \end{aligned}
    \end{equation}
\end{proposition}
\begin{proof}
    Suppose that \( (v^*, u^*) \) is an optimal solution \eqref{convex-forms:eq:unreduced-unconstrained-relu-mlp}.
    Defining \( r^* = v^* - u^* \), it holds by the triangle inequality that
    \begin{equation*}
        \norm{r_i^*}_2 \leq \norm{v_i^*}_2 + \norm{u_i^*}_2,
    \end{equation*}
    for each \( i \in [P] \).
    Since replacing \( v_i^* - u_i^* \) with \( r_i^* \) does not change the model prediction
    \[ \hat y = \sum_{D_i \in \tilde \calD} D_i X(v_i - u_i), \]
    we have shown that \( (r^*, 0) \) defines an equivalent model with the same or smaller objective value.
    Accordingly, any solution to~\eqref{convex-forms:eq:unreduced-unconstrained-relu-mlp} must have
    \( v_i^* = 0 \) or \( u_i^* = 0 \).
    Equivalence of the two problems follows immediately.
\end{proof}

\subsection{Extension to Multi-class Classification}\label{app:multi-class-extension}

Now we extend our sub-sampled C-ReLU and C-GReLU formulations to vector-valued
problems, such as occur in multi-class classification.
Our starting place is the following vector-output variant of the NC-ReLU problem
\begin{equation}\label{convex-forms:eq:vector-non-convex-relu-mlp}
    \min_{W^{(1)}, W^{(2)}} \mathcal{L}\Big(\sum_{i=1}^m
    (XW^{(1)}_{\cdot, i})_+ \sbr{W^{(2)}_{i}}^\top, Y\Big) +
    \frac{\lambda}{2}\sum_{i=1}^m \norm{W^{*(1)}_{\cdot, i}}_2^2 + \norm{W^{(2)}_{i}}_1^2,
\end{equation}
where now labels \(Y \in \R^{n \times C}\).
We note that the main difference between this formulation and
\cref{convex-forms:eq:non-convex-relu-mlp} is that each row of \(X\) now maps
to a vector rather than a single scalar, and the use of \(\ell_1\)-squared
regularization on the second-layer weights.
We now present a similar result to
\cref{lemma:mi-reformulation} for this particular problem.
\begin{lemma}
    The non-convex problem \eqref{convex-forms:eq:vector-non-convex-relu-mlp} is equivalent to the following program,
    \begin{equation}\label{convex-forms:eq:vector-mi-reformulation}
        \begin{aligned}
            \min_{\{W_{1}^k, w_{2}^k\}_{k=1}^C}  \; & \mathcal{L}\Big(\sum_{i=1}^m (XW^{*(1)}_{\cdot, i})_+ \sbr{W^{(2)}_{i}}^\top, Y\Big) + \lambda \sum_{i = 1}^{m}  \norm{W^{*(1)}_{\cdot, i}}_2 \\
                                                    & \text{\emph{s.t.}} \;\|W^{(2)}_{i}\|_1 = 1~\forall i \in [m]
        \end{aligned}
    \end{equation}
\end{lemma}
\begin{proof}
    Follow from the proof of \cref{lemma:mi-reformulation} (Appendix \ref{app:convex-forms}),
    i.e. apply Young's inequality to achieve
    \begin{align*}
        p^* & =
        \min_{W^{(1)}, W^{(2)}} \mathcal{L}\Big(\sum_{i=1}^m (XW^{*(1)}_{\cdot, i})_+ \sbr{W^{(2)}_{i}}^\top, Y\Big) + \frac{\lambda}{2}\sum_{i=1}^m \norm{W^{*(1)}_{\cdot, i}}_2^2 + \norm{W^{(2)}_{i}}_1^2 \\
            & = \min_{W^{(1)}, W^{(2)}} \mathcal{L}\Big(\sum_{i=1}^m (XW^{*(1)}_{\cdot, i})_+ \sbr{W^{(2)}_{i}}^\top, Y\Big) + \lambda\sum_{i=1}^m \norm{W^{*(1)}_{\cdot, i}}_2\norm{W^{(2)}_{i}}_1
    \end{align*}
    Then, this is clearly equivalent to
    \begin{equation*}
        \begin{aligned}
            \min_{W_{1}, W^{(2)}}  \; & \mathcal{L}\Big(\sum_{i=1}^m (XW^{*(1)}_{\cdot, i})_+ \sbr{W^{(2)}_{i}}^\top, Y\Big) + \lambda \sum_{i = 1}^{m}  \norm{W^{*(1)}_{\cdot, i}}_2 \\
                                      & \text{\emph{s.t.}} \;\|W^{(2)}_{i}\|_1 = 1~\forall i \in [m]
        \end{aligned}
    \end{equation*}
\end{proof}
We can form the one-vs-all convex reformulation as follows:
\begin{equation}\label{convex-forms:eq:vector-convex-relu-mlp}
    \begin{aligned}
        \min_{\{v^k, u^k\}_{k=1}^C} \, & \mathcal{L}\Big(\sum_{k=1}^C \sum_{D_i \in \tilde \calD} D_i X (v_{i}^k - u_{i}^k)e_k^\top, Y\Big)+\lambda\sum_{k=1}^C\sum_{D_i \in \tilde \calD}\norm{v_{i}^k}_2+\norm{u_{i}^k}_2 \\
                                       & \text{s.t.} \, \rbr{2 D_i - I_n} X v_{i}^k \succeq 0, \, \rbr{2 D_i - I_n} X u_{i}^k \succeq 0,
    \end{aligned}
\end{equation}
where \(e_k\) is the \(k\)th standard basis vector.

Then, we have the following analog of \cref{thm:duality-free} for the vector-output case:
\begin{restatable}{theorem}{convexVectorDualityFree}\label{thm:vector-duality-free}
    Suppose \(\rbr{W^{*(1)} W^{*(2)}}\) and \( \{\rbr{v^{k*}, u^{k*}}\}_{k=1}^C \) are global minimizers of Problems~\ref{convex-forms:eq:vector-mi-reformulation} and Problem~\ref{convex-forms:eq:vector-convex-relu-mlp}, respectively.
    If the number of hidden units satisfies
    \[
        m^* \geq b
        := \sum_{k=1}^C \sum_{D_i \in \tilde \calD} \abs{\cbr{v^{k*}_i : v^{k*}_i \neq 0}
            \cup \cbr{u^{k*}_i : u^{k*}_i \neq 0} },
    \]
    and the optimal activations are in the convex model,
    \[ \cbr{\text{\emph{diag}}\rbr{X W^{*(1)}_{\cdot, i} > 0 : i \in [m]} } \subseteq \tilde \calD, \]
    then the two problems have same the optimal value.
\end{restatable}
\begin{proof}
    We follow from the proof of \cref{thm:duality-free}.   Let \(p^*\) be the optimal value of \eqref{convex-forms:eq:vector-mi-reformulation} and \(d^*\) be the optimal value of \eqref{convex-forms:eq:vector-convex-relu-mlp}.

    First, suppose \( \{\rbr{v^{k*}, u^{k*}}\}_{k=1}^C\) is a global minimizer of Problem \ref{convex-forms:eq:vector-convex-relu-mlp}. Then, let
    \[ \rbr{W_{1(i, k)}^*, W_{2(i, k)}^*} = \bigcup_{D_i \in \tilde \calD} \bigcup_{k=1}^C \cbr{(v_i^{k*}, e_k): v_i^{k*} \neq 0} \cup \cbr{(u_i^{k*}, -e_k): u_i^{k*} \neq 0)}  \]
    where we set \(W_{1(i, k)}^* = 0\) and \(W_{2(i, k)}^*\) = 0 for non-assigned neurons. It holds by assumption that \(b \leq m\) and thus this is a valid input for \eqref{convex-forms:eq:vector-mi-reformulation}.
    Further, we have, due to the constraints,
    \begin{equation*}
        (XW_{1(i, k)}^*)_+W_{2(i, k)}^{*\top } = \begin{cases}
            D_iXv_{i}^{k*}e_k^\top & \text{if } W_{1(i, k)}^* = v_i^{k*} \\
            D_iXu_{i}^{k*}e_k^\top & \text{if } W_{1(i, k)}^* = u_i^{k*} \\
            0                      & o.w.
        \end{cases}
    \end{equation*}
    Inserting to the objective for the convex program,
    \begin{align}
        d^* & =\mathcal{L}\Big(\sum_{k=1}^C \sum_{D_i \in \tilde \calD} D_i X (v_{i}^{k*} - w_{i}^{k*})e_k^\top, Y\Big)+\lambda\sum_{k=1}^C\sum_{D_i \in \tilde \calD}\norm{v_{i}^{k*}}_2+\norm{w_{i}^{k*}}_2 \\
            & = \mathcal{L}\Big(\sum_{(i, k)}(XW_{1(i, k)}^*)_+W_{2(i, k)}^{*\top }, Y\Big)+\lambda\sum_{(i,k) }\norm{W_{1(i, k)}^*}_2                                                                        \\  &\geq p^*
    \end{align}
    Now, we seek to find the other direction, i.e. show \(p^* \geq d^*\) and show a mapping. Let \(\rbr{W^{*(1)}_{\cdot, i}, W^{*(2)}_{i}}\) be a solution to \eqref{convex-forms:eq:vector-mi-reformulation}. Defining, as in \cref{thm:duality-free},
    \[
        T(j) = \cbr{ i \in [m] : \text{diag}\rbr{X W^{*(1)}_{\cdot, i} > 0} = D_j }.
    \]
    Recalling
    \( \cbr{\text{diag}\rbr{X W^{*(1)}_{\cdot, i} > 0 : i \in [m]} } \subseteq \tilde \calD, \)
    by assumption, we define a valid candidate solution as
    \begin{align*}
        \cbr{\cbr{\rbr{v_j^{k*}, u_j^{k*}}}_{k \in [C]}}_{D_j \in \tilde \calD}
        &= \Bigg\{ \cbr{\sum_{i \in T(j)} W^{*(1)}_{\cdot, i} W^{*(2)}_{i, k} \mathds{1}(W^{(2)}_{i, k} \geq 0)}_{k \in [C]},\\
        & \hspace{4em} \cbr{-\sum_{i \in T(j)} W^{*(1)}_{\cdot, i} W^{*(2)}_{i, k} \mathds{1}(W^{(2)}_{i, k} < 0)}_{k \in [C]} \Bigg\}_{D_j \in \tilde \calD}
    \end{align*}
    Then, by the same co-linearity arguments as in \cref{thm:duality-free}, we have
    \begin{align*}
        p^* & =  \mathcal{L}\Big(\sum_{i=1}^m (XW^{*(1)}_{\cdot, i})_+ \sbr{W^{*(2)}_{i}}^{\top}, Y\Big) + \lambda \sum_{i = 1}^{m}  \norm{W^{*(1)}_{\cdot, i}}_2                                                        \\
            & = \mathcal{L}\Big(\sum_{i=1}^m \sum_{k=1}^C (XW^{*(1)}_{\cdot, i})_+ W^{(2)}_{i, k}e_k^\top, Y\Big) + \lambda \sum_{i = 1}^{m}  \norm{W^{*(1)}_{\cdot, i}}_2 \sum_{k=1}^C |W^{(2)}_{i, k*}|       \\
            & = \mathcal{L}\Big(\sum_{i=1}^m \sum_{k=1}^C (XW^{*(1)}_{\cdot, i})_+ W^{(2)}_{i, k}e_k^\top, Y\Big) + \lambda \sum_{i = 1}^{m} \sum_{k=1}^C \norm{W^{*(1)}_{\cdot, i}}_2  |W^{(2)}_{i, k}|       \\
            & = \mathcal{L}\Big(\sum_{D_j \in \tilde \calD} \sum_{k=1}^C D_jX(v_{j}^{k*}- w_{j}^{k*})e_k^\top, Y\Big) + \lambda \sum_{D_j \in \tilde \calD} \sum_{k=1}^C \norm{v_{j}^{k*}}_2 + \norm{w_{j}^{k*}}_2 \\
            & \geq d^*
    \end{align*}
\end{proof}
Thus, the vector-output NC-ReLU training problem~\eqref{convex-forms:eq:vector-non-convex-relu-mlp} is equivalent to the one-vs-all C-ReLU problem~\eqref{convex-forms:eq:vector-convex-relu-mlp} if the conditions of \cref{thm:vector-duality-free} are satisfied.
Further, taking \(\tilde \calD = \calD_X\) and applying \cref{thm:vector-duality-free} yields
\begin{equation*}
    {m}^{*} = \sum_{k=1}^C \sum_{D_i \in \calD_X} \abs{\cbr{v^{k*}_i : v^{k*}_i \neq 0} \cup \cbr{u^{k*}_i : u^{k*}_i \neq 0} }.
\end{equation*}
It follows that global optimization of the vector-output NC-ReLU problem requires \(m \geq m^*\) neurons, where \(m^* \leq C(n+1)\).

The gated ReLU analogs to vector-output ReLU architectures can be formulated in the same fashion.
